% Round 15. Multiplicity/native/pure-operational extension of Round 7.
% It identifies an obstruction to a supplied numerical model witness,
% not an obstruction whose nonvanishing disproves strong G.
\section{The canonical primitive residual on a fixed SNC model}
\label{r15res:section}

The component condition $\mathrm{LC}_2$ is sufficient because it
removes every numerical primitive codimension-two direction.
Without that condition, the missing directions can still be
identified exactly.  The result below gives a canonical residual
for each geometrically homologically trivial class on a fixed
model and with a fixed ample class.  Its vanishing is equivalent
to the existence of the stated numerical model lift; its
nonvanishing is not asserted to imply a nonzero actual
Bott--Chern obstruction or a failure of strong G.

\paragraph{Historical provenance and the exact extension.}
For a reduced strictly semistable special fiber, the compressed
divisor kernel, primitive normalization, canonical residual,
total isotropy, and the resulting positivity criteria were
already established in the frozen Round 7 sources
\path{history_01_07/bb_round7/proofs/threefold_codimension_two_obstruction.tex}
(\texttt{bb7tf:weighted-kernel},
\texttt{bb7tf:primitive-normalization},
\texttt{bb7tf:primitive-zero-energy}) and
\path{history_01_07/bb_round7/proofs/threefold_boundary_primitive_obstruction.tex}
(\texttt{bb7tf:canonical-primitive-obstruction},
\texttt{bb7tf:obstruction-totally-isotropic}).
Those statements already assume \emph{geometric} rational
$\ell$-adic triviality on the generic fiber; no weaker generic
cohomological hypothesis is claimed here.
The present section extends their argument to an SNC reduced
support with arbitrary positive component multiplicities,
retaining native degrees, and uses the final Round 12
pure-operational $K_0$ interface to certify the actual
strong-G conclusion.  The exact transpose-kernel quotient
below is an explicit linear-algebra consequence of the
compressed-kernel identity, not a claimed new independent
mechanism.  No novelty claim is made for the reduced case.

Use the hypotheses and notation of
Section~\ref{r15model:section}, except that
$\mathrm{LC}_2$ is \emph{not} assumed.  Thus the regular flat
projective model has relative dimension $n\geq3$, special
fiber $\sum_i m_iD_i$ with strict normal crossings reduced
support, and smooth closed strata over the finite residue
field.  The same relatively ample line $H$ is used on all
strata.  Put
\[
 V^a=\bigoplus_i N^a(D_i)_{\mathbf Q},\qquad
 \Lambda=L^{n-3},\qquad
 P_H^2=\ker(\Lambda:V^2\longrightarrow V^{n-1}),
\tag{NR15.1}\label{r15res:primitive-space}
\]
where $L$ means intersection with $H$.  The same symbol
$\Lambda$ is used in its two degrees
$V^1\to V^{n-2}$ and $V^2\to V^{n-1}$.
Numerical intersection gives the perfect rational pairings
between complementary $V^a$.  All degrees are native degrees
over the original residue field.  The maps
$A_a:V^a\to V^{a+1}$ are the actual vertical incidence maps
of \eqref{r15model:A}, without a change of multiplicity convention.

\begin{lemma}[Divisor Lefschetz linear algebra without $\mathrm{LC}_2$]
\label{r15res:divisor-linear-algebra}
The map $L^{n-2}:V^1\to V^{n-1}$ is an isomorphism.
Consequently $\Lambda:V^1\to V^{n-2}$ is injective,
$\Lambda:V^2\to V^{n-1}$ is surjective, and
\[
                    V^2=L V^1\oplus P_H^2.
\tag{NR15.2}\label{r15res:lefschetz-splitting}
\]
The perfect pairing between $V^2$ and $V^{n-2}$ induces a
perfect pairing
\[
            P_H^2\ \times\ (V^{n-2}/\Lambda V^1)
                                          \longrightarrow\mathbf Q.
\tag{NR15.3}\label{r15res:primitive-duality}
\]
These assertions use only the numerical divisor Hodge index theorem.
\end{lemma}

\begin{proof}
Component by component, the divisor pairing
$\deg(u v H^{n-2})$ is nondegenerate.  It makes
$L^{n-2}:N^1(D_i)\to N^{n-1}(D_i)$ injective, and the two
spaces have the same finite dimension by numerical duality.
It is therefore an isomorphism.
Since $L^{n-2}=L\Lambda$ on $V^1$, the first $\Lambda$ is
injective.  Since $L^{n-2}=\Lambda L$ there, the second
$\Lambda$ is surjective.  The latter is an isomorphism on
$L V^1$, proving \eqref{r15res:lefschetz-splitting}.
The annihilator of $\Lambda V^1$ in $V^2$ is exactly
$\ker\Lambda=P_H^2$, by associativity of intersection.
Perfect numerical duality proves \eqref{r15res:primitive-duality}.
No cycle-class map on $N^2$ and no codimension-two sign
conjecture are used.
\end{proof}

\begin{lemma}[Compression does not enlarge the divisor kernel]
\label{r15res:compressed-kernel}
For every model under consideration, without $\mathrm{LC}_2$,
\[
              \boxed{\ \ker(\Lambda A_1)=\ker A_1.\ }
\tag{NR15.4}\label{r15res:kernel-compression}
\]
Equivalently,
\[
                       \operatorname{im}A_1\cap P_H^2=0.
\tag{NR15.5}\label{r15res:image-primitive-disjoint}
\]
\end{lemma}

\begin{proof}
We indicate exactly how the frozen divisor-energy proof gives
the stronger statement.  Let $M$ multiply each component by
$m_i$, let $\widetilde A=M A_1 M$, and let
$\gamma_0^m=M^{-1}\gamma_0D_c$ be the weighted edge Gysin.
Use the edge functional denoted $\lambda(w)$ in Round 12;
to avoid confusing it with $L$, denote it here by $\ell(w)$.
Thus on a double stratum $D_e=D_i\cap D_j$,
\[
       \ell(w)_e=m_im_j\deg_{D_e}((w_j-w_i)|_{D_e}H^{n-2}).
\tag{NR15.6}\label{r15res:edge-functional}
\]
The frozen source
\path{round12/numerical_lift/curve_lifts_and_actual_green.tex}
proves the following identities before its final kernel
conclusion: (NL11), (NL12), (NL13), and (NL14).  In the present
notation they give
\[
 \begin{split}
 \ell(\gamma_0^m h)&=-d_1^tD_b d_1h,\\
 (d_0^t\ell(w))_i&=-\deg_{D_i}((\widetilde Aw)_iH^{n-2}),\\
 \ell(w)^th&=\sum_i m_i\deg_{D_i}
                            (w_i(\gamma_0^m h)_iH^{n-2}),\\
 \gamma_0^m(\ker d_1\cap\ker(d_0^tD_a))&=0.
 \end{split}
\tag{NR15.7}\label{r15res:frozen-identities}
\]
Here $D_a,D_b$ have the positive native edge and triple-stratum
weights of that source.  In particular the last identity is
independent of any assumption on $w$.  Also
$\widetilde A\gamma_0^m=0$ on these numerical groups.

Suppose only that $\Lambda\widetilde Aw=0$.
The second identity in \eqref{r15res:frozen-identities}
still gives $d_0^t\ell(w)=0$, since its right side is the
degree of $H\Lambda\widetilde Aw$.
The last two identities make $\ell(w)$ vanish on
$\ker d_1\cap\ker(d_0^tD_a)$.
The positive-weight decomposition
\[
 \ker d_1=\operatorname{im}d_0\oplus
                  (\ker d_1\cap\ker(d_0^tD_a))
\]
therefore implies $\ell(w)|_{\ker d_1}=0$.
As in the frozen rational linear algebra, there is a rational
edge vector $a$ with
$d_1^tD_b d_1a=\ell(w)$.
Set $w_0=w+\gamma_0^m a$.  Then
\[
               \ell(w_0)=0,
               \qquad\Lambda\widetilde Aw_0=0.
\tag{NR15.8}\label{r15res:adjusted-weak-kernel}
\]

The weighted second-energy identity (NL14) now reads
\[
 \begin{split}
 &\sum_{e=(i,j)}m_im_j
             \deg_{D_e}\bigl(((w_{0,j}-w_{0,i})|_{D_e})^2H^{n-3}\bigr)\\
 &\hspace{15mm}=-\sum_i
             \deg_{D_i}(w_{0,i}(\widetilde Aw_0)_iH^{n-3})=0.
 \end{split}
\tag{NR15.9}\label{r15res:weak-energy}
\]
The final equality follows from the second assertion of
\eqref{r15res:adjusted-weak-kernel}; it does not require
$\widetilde Aw_0=0$ in advance.
The first equality is a termwise identity for arbitrary $w_0$:
grouping the two endpoint contributions of each edge in (NL6)
gives the negative of its weighted squared difference.
Thus that equality has no hidden full-kernel hypothesis.
The first assertion makes every divisor difference on a
double stratum primitive for $H$.  Divisor Hodge index on
those smooth projective $(n-1)$-folds makes every summand
nonpositive, with equality only for a numerically zero
divisor.  All $m_im_j$ are positive, so every difference
vanishes numerically.  The explicit formula (NL6) for
$\widetilde A$ then gives $\widetilde Aw_0=0$.
As $\widetilde A\gamma_0^m=0$, also $\widetilde Aw=0$.

For an original vector $u$ with $\Lambda A_1u=0$, put
$w=M^{-1}u$.  The multiplication matrix $M$ commutes with
$\Lambda$, so the argument applies and gives $A_1u=0$.
The reverse inclusion is immediate.  Formula
\eqref{r15res:image-primitive-disjoint} is the same assertion
on the image.
\end{proof}

\begin{proposition}[The exact residual kernel quotient]
\label{r15res:kernel-quotient}
There is an exact sequence of rational vector spaces
\[
 0\longrightarrow\Lambda\ker A_1
   \longrightarrow\ker A_{n-2}
   \longrightarrow V^{n-2}/\Lambda V^1
   \longrightarrow0,
\tag{NR15.10}\label{r15res:exact-sequence}
\]
where the last map is the natural quotient map.
Consequently
\[
 \frac{\ker A_{n-2}}{\Lambda\ker A_1}
          \simeq V^{n-2}/\Lambda V^1,
 \qquad
 \left(\frac{\ker A_{n-2}}{\Lambda\ker A_1}\right)^*
          \simeq P_H^2.
\tag{NR15.11}\label{r15res:canonical-residual-space}
\]
\end{proposition}

\begin{proof}
Restriction, Gysin, and the self-normal term commute with $H$,
so $A_{n-2}\Lambda=\Lambda A_1$.
Together with Lemma~\ref{r15res:compressed-kernel}, this gives
\[
       \ker A_{n-2}\cap\Lambda V^1=\Lambda\ker A_1.
\]
It proves exactness at the first two terms and injectivity
of the indicated quotient map.

The numerical transpose of $A_1$ is $A_{n-2}$, by the
restriction--Gysin projection formula and the symmetric normal
self-intersection term.  Put $d_1=\dim V^1$,
$d_2=\dim V^2=\dim V^{n-2}$, and $a=\operatorname{rank}A_1$.
The source quotient in \eqref{r15res:canonical-residual-space}
has dimension
$(d_2-a)-(d_1-a)=d_2-d_1$.
The target quotient has that same dimension because
$\Lambda:V^1\to V^{n-2}$ is injective.
Thus the injective quotient map is surjective as well.
The dual identification is \eqref{r15res:primitive-duality}.
\end{proof}

\begin{theorem}[Canonical primitive component residual]
\label{r15res:canonical-residual}
Let $z\in\operatorname{CH}^2(X)_{\mathbf Q}$ have zero
geometric rational $\ell$-adic class for one $\ell\ne p$.
For an arbitrary Chow extension $\Theta$ to $\mathcal X$, put
\[
          T=(j_i^!\Theta)_i\in V^2,
          \qquad B_H=\Lambda A_1:V^1\longrightarrow V^{n-1}.
\tag{NR15.12}\label{r15res:compressed-system}
\]
The system
\[
                         B_Hu=\Lambda T
\tag{NR15.13}\label{r15res:normalization-equation}
\]
always has a rational solution.  The residual
\[
              R_{\mathcal X,H}(z)=T-A_1u\in P_H^2
\tag{NR15.14}\label{r15res:residual}
\]
is independent of that solution and of the Chow extension
$\Theta$.  It depends only on the class $z$, the supplied
model, and the chosen ample class, and is rational-linear in $z$.
There are equivalent conditions:
\begin{enumerate}
\item $R_{\mathcal X,H}(z)=0$;
\item some, equivalently any, extension of $z$ can be corrected
by a sum of vertical component divisors so that all component
restrictions are zero in $N^2(D_i)$;
\item the functional $t_\Theta(v)=\sum_i\deg(T_i v_i)$
vanishes on $\ker A_{n-2}$.
\end{enumerate}
The common zero condition is independent of the choice of $H$.
When it holds, the frozen pure-operational model-Chern detector
and exact-support bridge give strong G for every specified
cycle representing $z$, with all original actual tests.
No converse from actual obstruction vanishing or from strong G
to these numerical conditions is asserted.
\end{theorem}

\begin{proof}
We first recall precisely the annihilation that does not need
$\mathrm{LC}_2$.  The frozen divisor kernel is
$\ker A_1=M\ker\delta+\operatorname{im}\gamma_0$.
For a signed edge Gysin $\gamma_0a$, the two copies of its
double-stratum cycle cancel after pushforward to $\mathcal X$.
Restricting $\Theta$ to that stratum and intersecting with
$H^{n-3}$ therefore gives
$t_\Theta(\Lambda\gamma_0a)=0$ using supported degrees.
For a compatible divisor tuple $w\in\ker\delta$, the frozen
cohomological-lift theorem \texttt{r12n:cohom-lift} and
multiplicity-weighted trace theorem
\texttt{r12n:weighted-trace}, applied to
\[
       \operatorname{cl}(\Theta)\cup b\cup c_1(H)^{n-3},
\]
give $t_\Theta(\Lambda Mw)=0$, since its generic restriction
contains $\operatorname{cl}_\ell(z)=0$.
The cohomology class of $\Theta$ is taken via its pure
operational $K_0$ representative, as in
Theorem~\ref{r15model:lift}; no cycle-class map out of $N^2$
is assumed.  These two arguments show
\[
                     t_\Theta|_{\Lambda\ker A_1}=0.
\tag{NR15.15}\label{r15res:known-annihilation}
\]

The operator $B_H$ is symmetric under the perfect pairing
of $V^1$ with $V^{n-1}$.  Indeed its transpose is
$A_{n-2}\Lambda=\Lambda A_1=B_H$.
By Lemma~\ref{r15res:compressed-kernel},
$\ker B_H=\ker A_1$.
Equation \eqref{r15res:known-annihilation} says exactly that
$\Lambda T$ annihilates this kernel.  Finite-dimensional
rational linear algebra for the symmetric map $B_H$ therefore
puts $\Lambda T$ in its image.  This proves existence in
\eqref{r15res:normalization-equation}.

That equation makes $R=T-A_1u$ lie in $P_H^2$.
Two solutions differ by $\ker B_H=\ker A_1$, so they give
the same $R$.  If $\Theta'$ is another Chow extension,
closed--open Chow localization expresses
\[
                   \Theta'-\Theta=\sum_i j_{i*}v_i,
                       \qquad v_i\in\operatorname{CH}^1(D_i)_{\mathbf Q}.
\tag{NR15.16}\label{r15res:localization}
\]
For completeness, localization gives a cycle supported on
the reduced special fiber; every codimension-two component
of that supported cycle is a divisor in at least one $D_i$,
so its cycle lifts to the displayed sum.  The scheme
multiplicities $m_i$ are not inserted in this localization
map.  They already occur in $A_1$ and in the compatible
divisor kernel.  This is Chow localization over an excellent
DVR with its dimension function, for example Stacks Project
Tag~02RX, Lemma~42.19.3, applied in codimension two.
It does not use an unweighted degree map on the model's
zero-cycles.

Numerically, \eqref{r15res:localization} gives
$T'=T+A_1v$.  If $u$ solves the old system, $u+v$ solves
the new one and gives the identical residual.  This proves
extension independence and also independence of the
representative of the generic Chow class.  Linearity follows
from linearity of the equations and uniqueness of the residual.

If $R=0$, a chosen rational divisor representative for each
$u_i$ corrects $\Theta$ to an extension with all numerical
component restrictions zero.  Conversely, such a correction
gives a vector $u$ with $A_1u=T$, hence solves
\eqref{r15res:normalization-equation} and makes $R=0$.
The equivalence with the third condition is the perfect
transpose identity
$\operatorname{im}A_1=(\ker A_{n-2})^\perp$.
Alternatively, for $v\in\ker A_{n-2}$ one has
\[
       t_\Theta(v)=\langle A_1u+R,v\rangle=\langle R,v\rangle,
\tag{NR15.17}\label{r15res:functional-realization}
\]
and \eqref{r15res:canonical-residual-space} makes this
functional zero exactly when $R=0$.
The numerical correction condition does not mention $H$,
so its equivalent zero condition is independent of that choice.

Finally the corrected class has the pure operational virtual
bundle supplied by frozen \texttt{r12n:K0} in degree two.
Its supported degree on every integral special-fiber surface
is zero, because that surface lies in some $D_i$ and its
pairing with the corresponding $N^2$ class is zero.
The frozen detector and exact-support bridge then apply
exactly as in Theorem~\ref{r15model:lift}.  They use the
established sufficient direction only.
\end{proof}

\begin{corollary}[The precise remaining numerical directions]
\label{r15res:interpretation}
For $n=3$ the primitive residual space is zero.  For arbitrary
$n\geq3$, condition $\mathrm{LC}_2$ is equivalent to
$P_H^2=0$, so it forces the residual to vanish for every
geometrically homologically trivial codimension-two class.
Without that condition, every such class still admits a
vertical divisor correction whose component restrictions are
the unique primitive tuple \eqref{r15res:residual}.
\end{corollary}

\begin{proof}
For $n=3$, $\Lambda$ is the identity on $V^2$.
In every dimension, \eqref{r15res:lefschetz-splitting} identifies
surjectivity of $L:V^1\to V^2$ with $P_H^2=0$, component
by component.  The remaining assertion is the defining
normalization equation of the theorem.
\end{proof}

\paragraph{Open mathematical issue.}
No nonzero value of $R_{\mathcal X,H}$ for a homologically
trivial generic class is constructed here, and no universal
vanishing theorem for it is proved.  Its codomain is a
numerical primitive space on the special components, not an
actual Aeppli or Bott--Chern group.  A nonzero value would
exclude this numerical lift on this model; it would not
exclude a different model witness, another actual primitive,
or a strong Green form.  Determining whether and how generic
homological triviality controls these remaining primitive
directions is the precise next numerical problem.
